\section{Methods}
\label{sec:methods}

\subsection{Instrument and firmware baseline}
The instrument under test is HunStat2, a single-channel three-electrode
potentiostat built around the AD5941 AFE and a Seeed Studio XIAO RP2040
microcontroller, following the AD5940 datasheet, standard three-electrode
potentiostat design practice, and the AD5940 evaluation-kit reference
firmware~\cite{devices2024ad5940,ad5940_eis_github}, detailed in
Section~\ref{sec:instrument-design}. The firmware baseline for this
study is the \texttt{AD5941\_25} sketch as of the ``update~7 (20 September
26)'' snapshot in the project repository, compiled with the
\texttt{rp2040:rp2040} Arduino core (v5.6.0) via \texttt{arduino-cli}, and
flashed to a physical XIAO RP2040 board over USB for every measurement
reported here. No simulation or mock serial layer was used at any point;
every dataset in this paper was collected from the physical instrument.

\subsection{Software-architecture audit}
Every \texttt{.h}/\texttt{.cpp} file under \texttt{AD5941\_25/src/} was read
in full and characterized on five axes: what class or module owns the file,
what OOP concepts (encapsulation, inheritance, polymorphism) are actually
present versus merely named, what the file depends on, what problems the
implementation has, and what a targeted fix would look like. A repository-wide
search for \texttt{virtual}, \texttt{: public} (inheritance), and
\texttt{template<} was run to establish, independent of any per-file
judgment, whether polymorphism or generic code exists anywhere in project
code. Call-site tracing -- following every function actually invoked from
\texttt{loop()} down through the dispatch layer -- was used to distinguish
code that executes from code that merely compiles.

\subsection{Refactor}
Three changes were made, each scoped to be individually verifiable against
the pre-refactor source:

\begin{enumerate}
\item \textbf{Shared base class for DC-excitation methods.} The audit found
that \texttt{C\_CA::ConfigDCMeasurement()}, \texttt{C\_DPV::ConfigDCMeasurement()},
and \texttt{C\_SWV::ConfigDCMeasurement()} -- along with each class's
\texttt{MeasureCurrentRaw()} and \texttt{RawToCurrent()} -- were identical
except for class-qualification and comment wording, confirmed by a line-level
diff of all three implementations. These three methods configure the AD5940
low-power loop (LPDAC/LPTIA), trigger one ADC conversion, and convert the raw
code to amperes; they encode a substantial amount of hard-won,
previously-debugged register-sequencing knowledge (documented in the source
comments: AFE wake state, DAC polarity, PGA gain enumeration, LPTIA routing).
A new class, \texttt{C\_DCPotentiostatMethod}, was introduced holding these
three methods verbatim -- copied register-for-register and
formula-for-formula, not rewritten -- and \texttt{C\_CA}, \texttt{C\_DPV},
and \texttt{C\_SWV} were changed to inherit from it publicly, retaining only
their method-specific \texttt{Run()} sequencing (the differing pulse/staircase
timing that makes CA, DPV, and SWV distinct methods). This is a standard
Extract Superclass refactor~\cite{fowler2018refactoring}: it is a template
method in spirit (a shared measurement primitive, per-method sequencing on
top) without requiring a \texttt{virtual} dispatch, since the concrete
subclass is always known statically at the call site.
\item \textbf{Dead-code removal.} \texttt{command\_processing.cpp/.h} (a
$\sim$700-line free-function command parser) and
\texttt{electrochemical\_methods.cpp/.h} (a free-function
reimplementation of CA/SWV/DPV against a separate set of globals) were
confirmed, by tracing every call site, to be unreachable: \texttt{loop()}
calls only \texttt{C\_Communication::ReadAndProcess()}, whose own
\texttt{SplitAndProcessCommands()} is a same-named but distinct class
method, never the free function of the same name. Both files were deleted
rather than merely un-included, since the Arduino build compiles every
\texttt{.cpp} file under the sketch's \texttt{src/} tree regardless of
whether it is \texttt{\#include}d.
\item \textbf{Status-LED bug fix.} \texttt{Interface\_SetLed(color)}
unconditionally rendered white regardless of its argument, silently
defeating every caller's status color coding (cyan during CA, magenta
during DPV, yellow during SWV, green on completion). The correct
color-index-to-RGB mapping was recovered from an already-correct but
separately dead helper (\texttt{LED()} in \texttt{utilities.cpp}, never
called from anywhere) and used to fix the live function; the previously
dead helper was left untouched and out of scope.
\end{enumerate}

Code duplication of this kind is not a cosmetic concern: independent
empirical work on code clones has shown that inconsistently-maintained
duplicates measurably increase defect rates relative to non-duplicated
code~\cite{juergens2009clones}, which is exactly the failure mode two
divergent globals-based implementations of the same command grammar
represent.

\subsection{Dummy-cell validation protocol}
\label{sec:protocol}
Correctness preservation was checked empirically rather than argued from
the diff alone, reusing the dummy-cell protocol validated in the companion
manuscript~\cite{clement2026linearity}. Two real loads were used, connected
in place of an electrochemical cell across the working, reference, and
counter electrode terminals:

\begin{itemize}
\item \textbf{Randles dummy cell} ($R_s=\SI{560}{\ohm}$,
  $R_{ct}=\SI{10}{\kilo\ohm}$, $C_{dl}=\SI{33}{\nano\farad}$), used for the
  primary before/after comparison.
\item \textbf{Plain resistor} ($\SI{10}{\kilo\ohm}$, no capacitive branch),
  substituted in place of the Randles cell as a second, independent load
  used only on the refactored firmware, to corroborate that the measured
  current-voltage relationship is sane under a load the refactor was not
  implicitly tuned against.
\end{itemize}

For each firmware build (pre-refactor, post-refactor), the board was
flashed, and all four electrochemical methods were exercised over the
Randles cell using the instrument's native ASCII command protocol over USB
serial at \SI{1}{\mega baud}:

\begin{itemize}
\item \textbf{CV}: \texttt{D-200,600,5,100,1} (start \SI{-200}{\milli\volt},
  stop \SI{600}{\milli\volt}, step \SI{5}{\milli\volt}, scan rate
  \SI{100}{\milli\volt\per\second}, 1 cycle), triggered with \texttt{M}.
\item \textbf{CA}: step to \SI{100}{\milli\volt} for \SI{2}{\second} at
  \SI{40}{\hertz} sampling, triggered with \texttt{A}.
\item \textbf{SWV}: \SIrange{0}{200}{\milli\volt}, \SI{20}{\milli\volt}
  step, \SI{25}{\milli\volt} amplitude, triggered with \texttt{W}.
\item \textbf{DPV}: \SIrange{0}{200}{\milli\volt}, \SI{20}{\milli\volt}
  step, \SI{25}{\milli\volt} pulse amplitude, triggered with \texttt{D}.
\end{itemize}

Each sweep was captured over serial with a Python client enforcing an
explicit per-step deadline (never blocking indefinitely on a hung board),
parsing the method-tagged output lines (\texttt{CV,<mV>,<I>},
\texttt{CA,<s>,<I>,<raw>,<timeout flag>}, \texttt{SWV,<mV>,<$\Delta$I>},
\texttt{DPV,<mV>,<$\Delta$I>}) until the method's completion marker
(\texttt{*\_END}) or a 2-second data-silence timeout. The same procedure
was repeated on the refactored firmware against the $\SI{10}{\kilo\ohm}$
resistor load. No firmware parameter, calibration constant, or command
sequence differed between the pre- and post-refactor runs; the only
variable was the source code itself.
